% !TEX root = chapter-standalone.tex

\section{The case of semicategories}
\label{sec:parallel-for-semicategories}



\subsection{Stacking}

\begin{ctdefinition}[Stacking semicategory]
    \label{def:simple-stacking-semi-cat}
    A \maindef{stacking semicategory} is a \SY{semicategory}~\CatC with the following additional constituents and properties.

    \constit
    \begin{itemize}
        \item A stacking operation~$\mtimescatob \colon \ObC \cartprod \ObC \sto \ObC$.
        \item A stacking operation~$\mtimescatmor \colon \Mor_\CatC \cartprod \Mor_\CatC \sto \Mor_\CatC$.
    \end{itemize}

    \condit
    \begin{itemize}
        \item The two operations~$\mtimescatob$ and~$\mtimescatmor$ are compatible in the sense that
              %
              \begin{equation}
                  \prfperiod{
                      \moran{1} \colon \Objan{1} \mto \Objbn{1}
                      \qquad
                      \moran{2} \colon \Objan{2} \mto \Objbn{2}
                  }{
                      \moran{1} \mtimescatmor \moran{2} \colon  \Objan{1} \mtimescatob  \Objan{2} \mto \Objbn{1} \mtimescatob \Objbn{2}
                  }
              \end{equation}
    \end{itemize}
\end{ctdefinition}


\subsection{Functorial stacking}

\begin{ctdefinition}[Functorial stacking semicategory]
    \label{def:functorial-stacking-semi-cat}
    A \maindef{functorial stacking semicategory} is a stacking \SY{semicategory} where the two stacking operations~$\mtimescatob$ and~$\mtimescatmor$ are the two components of a \SY{functor} \begin{equation}
        \mtimescat\colon \CatC \Ctimes \CatC \to \CatC.
    \end{equation}
    In infix notation, this means that
    \begin{equation}
        (\mora \mthen \morc )
        \mtimescatmor (\morb \mthen \mord )
        =
        (\mora \mtimescatmor \morb )
        \mthen (\morc \mtimescatmor \mord )
    \end{equation}
    for all morphisms~$\mora, \morb, \morc$, and $\mord$ (where respectively $\mora$ and $\morc$, and $\morb$ and $\mord$ are composable), and that
        \begin{equation}
        \catidat{\Obja} \mtimescatmor \catidat{\Objb} = \catidat{\Obja \mtimescatob \Objb}
    \end{equation}
    for all objects $\Obja$, $\Objb$ of $\CatC$, whenever all three of these identity morphisms exist.
\end{ctdefinition}

\begin{example}
    We want to show that \Moore, equipped with the defined stacking operations, is almost a \SY{functorial stacking semicategory},
    but not quite.
    Consider four Moore machines:
    \begin{equation*}
        \mora\colon \prin_\mora \mto \prout_\mora,\quad \morb\colon \prin_\morb \mto \prout_\morb, \quad \morc\colon \prout_\mora \mto \prout_\morc,\quad \mord\colon \prout_\morb \mto \prout_\mord.
    \end{equation*}
    We want to check if the equation
    \begin{equation*}
        (\mora \mthen \morc)
        \mtimescatmor (\morb\mthen \mord)=(\mora\mtimescatmor \morb)\mthen (\morc\mtimescatmor \mord)
    \end{equation*}
    holds.
    Let's start from the left-hand side.
    First, we have:
    \begin{equation*}
        \begin{aligned}
            \mora \mthen \morc & =\tup{\prin_\mora, \prst_\mora \cprod \prst_\morc, \prout_\morc, \prdyn_{\mora\mthen \morc},\prreadout_{\mora\mthen \morc}, \prstart_\mora\tupconcat \prstart_\morc}, \\
            \morb \mthen \mord & =\tup{\prin_\morb, \prst_\morb \cprod \prst_\mord, \prout_\mord, \prdyn_{\morb\mthen \mord},\prreadout_{\morb\mthen \mord}, \prstart_\morb\tupconcat \prstart_\mord},
        \end{aligned}
    \end{equation*}
    with
    \begin{equation*}
        \defmapcomma{\prdyn_{\mora\mthen \morc}}
        {\prin_\mora\cprod \prst_\mora \cprod \prst_\morc}
        {\sto}
        {\prst_\mora \cprod \prst_\morc}
        {\prinel_\mora\tupconcat \prstel_\mora \tupconcat \prstel_\morc}
        {\prdyn_\mora(\prinel_\mora \tupconcat \prstel_\mora)\tupconcat \prdyn_\morc(\prreadout_\mora(\prstel_\mora)\tupconcat \prstel_\morc)}
    \end{equation*}
    %
    \begin{equation*}
        \defmapcomma{\prdyn_{\morb\mthen \mord}}
        {\prin_\morb\cprod \prst_\morb \cprod \prst_\mord}
        {\sto}
        {\prst_\morb \cprod \prst_\mord}
        {\prinel_\morb\tupconcat \prstel_\morb \tupconcat \prstel_\mord}
        {\prdyn_\morb(\prinel_\morb \tupconcat \prstel_\morb)\tupconcat \prdyn_\mord(\prreadout_\morb(\prstel_\morb)\tupconcat \prstel_\mord)}
    \end{equation*}
    %
    \begin{equation*}
        \defmapcomma{\prreadout_{\mora\mthen \morc}}
        {\prst_\mora \cprod \prst_\morc}
        {\sto}
        {\prout_\morc}
        {\prstel_\mora \tupconcat \prstel_\morc}
        {\prreadout_\morc(\prstel_\morc)}
    \end{equation*}
    and
    \begin{equation*}
        \defmapperiod{\prreadout_{\morb\mthen \mord}}
        {\prst_\morb \cprod \prst_\mord}
        {\sto}
        {\prout_\mord}
        {\prstel_\morb \tupconcat \prstel_\mord}
        {\prreadout_\mord(\prstel_\mord)}
    \end{equation*}

    Furthermore:
    \begin{widepar}
        \begin{equation*}
            (\mora \mthen \morc)\mtimescatmor (\morb \mthen \mord)=\tup{\prin_\mora\cprod \prin_\morb, \prst_\mora \cprod \prst_\morc\cprod \prst_\morb \cprod \prst_\mord, \prout_\morc\cprod \prout_\mord,
                \prdyn_{(\mora\mthen \morc)\mtimescatmor(\morb\mthen \mord)}, \prreadout_{(\mora\mthen \morc)\mtimescatmor(\morb\mthen \mord)}, \prstart_\mora\tupconcat \prstart_\morc\tupconcat \prstart_\morb\tupconcat \prstart_\mord},
        \end{equation*}
    \end{widepar}
    with
    \begin{widepar}
        \begin{equation*}
            \defmapcomma{\prdyn_{(\mora\mthen \morc)\mtimescatmor(\morb\mthen \mord)}}
            {\prin_\mora\cprod \prin_\morb\cprod \prst_\mora \cprod \prst_\morc\cprod \prst_\morb \cprod \prst_\mord}
            {\sto}
            {\prst_\morc\cprod \prst_\morb \cprod \prst_\mord}
            {\prinel_\mora\tupconcat \prinel_\morb \tupconcat \prstel_\mora \tupconcat \prstel_\morc \tupconcat \prstel_\morb \tupconcat \prstel_\mord}
            {\underbrace{\prdyn_{\mora\mthen \morc}(\prinel_\mora\tupconcat \prstel_\mora \tupconcat \prstel_\morc)\tupconcat \prdyn_{\morb\mthen \mord}(\prinel_\morb\tupconcat \prstel_\morb \tupconcat \prstel_\mord)}_{(1)}}
        \end{equation*}
    \end{widepar}
    where
    \begin{equation*}
        (1)=\prdyn_\mora(\prinel_\mora \tupconcat \prstel_\mora)\tupconcat \prdyn_\morc(\prreadout_\mora(\prstel_\mora)\tupconcat \prstel_\morc) \tupconcat
        \prdyn_\morb(\prinel_\morb \tupconcat \prstel_\morb)\tupconcat \prdyn_\mord(\prreadout_\morb(\prstel_\morb)\tupconcat \prstel_\mord),
    \end{equation*}
    and
    \begin{equation*}
        \defmapcomma{\prreadout_{(\mora\mthen \morc)\mtimescatmor(\morb\mthen \mord)}}
        {\prst_\mora \cprod \prst_\morc\cprod \prst_\morb \cprod \prst_\mord}
        {\sto}
        {\prout_\morc\cprod \prout_\mord}
        {\prstel_\mora \tupconcat \prstel_\morc \tupconcat \prstel_\morb \tupconcat \prstel_\mord}
        {\underbrace{\prreadout_{\mora\mthen \morc}(\prstel_\mora \tupconcat \prstel_\morc)\tupconcat \prreadout_{\morb\mthen \mord}(\prstel_\morb \tupconcat \prstel_\mord)}_{(2)}}
    \end{equation*}
    where
    \begin{equation*}
        (2)=\prreadout_\morc(\prstel_\morc)\tupconcat \prreadout_\mord(\prstel_\mord).
    \end{equation*}
    On the other hand, we have:
    \begin{equation*}
        \begin{aligned}
            \mora \mtimescatmor \morb & =\tup{\prin_\mora \cprod \prin_\morb, \prst_\mora \cprod \prst_\morb, \prout_\mora \cprod \prout_\morb, \prdyn_{\mora\mtimescatmor \morb},\prreadout_{\mora\mtimescatmor \morb}, \prstart_\mora\tupconcat \prstart_\morb}, \\
            \morc \mtimescatmor \mord & =\tup{\prin_\morc \cprod \prin_\mord, \prst_\morc \cprod \prst_\mord, \prout_\morc \cprod \prout_\mord, \prdyn_{\morc\mtimescatmor \mord},\prreadout_{\morc\mtimescatmor \mord}, \prstart_\morc\tupconcat \prstart_\mord},
        \end{aligned}
    \end{equation*}
    with
    \begin{equation*}
        \defmapcomma{\prdyn_{\mora \mtimescatmor \morb}}{\prin_\mora \cprod \prin_\morb \cprod \prst_\mora \cprod \prst_\morb}
        {\sto}{\prst_\mora \cprod \prst_\morb}
        {\prinel_\mora \tupconcat \prinel_\morb \tupconcat \prstel_\mora \tupconcat \prstel_\morb}
        {\prdyn_\mora(\prinel_\mora \tupconcat \prstel_\mora)\tupconcat \prdyn_\morb(\prinel_\morb \tupconcat \prstel_\morb)}
    \end{equation*}
    %
    \begin{equation*}
        \defmapcomma{\prdyn_{\morc \mtimescatmor \mord}}{\prin_\morc \cprod \prin_\mord \cprod \prst_\morc \cprod \prst_\mord}
        {\sto}{\prst_\morc \cprod \prst_\mord}
        {\prinel_\morc \tupconcat \prinel_\mord \tupconcat \prstel_\morc \tupconcat \prstel_\mord}
        {\prdyn_\morc(\prinel_\morc \tupconcat \prstel_\morc)\tupconcat \prdyn_\mord(\prinel_\mord \tupconcat \prstel_\mord)}
    \end{equation*}
    %
    \begin{equation*}
        \defmapcomma{\prreadout_{\mora \mtimescatmor \morb}}
        {\prst_\mora \cprod \prst_\morb}{\sto}{\prout_\mora \cprod \prout_\morb}
        {\prstel_\mora \tupconcat \prstel_\morb}{\prreadout_\mora(\prstel_\mora)\tupconcat \prreadout_\morb(\prstel_\morb)}
    \end{equation*}
    and
    \begin{equation*}
        \defmapperiod{\prreadout_{\morc \mtimescatmor \mord}}
        {\prst_\morc \cprod \prst_\mord}{\sto}{\prout_\morc \cprod \prout_\mord}
        {\prstel_\morc \tupconcat \prstel_\mord}{\prreadout_\morc(\prstel_\morc)\tupconcat \prreadout_\mord(\prstel_\mord)}
    \end{equation*}

    Furthermore:
    \begin{widepar}
        \begin{equation*}
            (\mora \mtimescatmor \morb)\mthen (\morc \mtimescatmor \mord)=\tup{\prin_\mora\cprod \prin_\morb, \prst_\mora \cprod \prst_\morb\cprod \prst_\morc \cprod \prst_\mord, \prout_\morc\cprod \prout_\mord,
                \prdyn_{(\mora\mtimescatmor \morb)\mtimescatmor(\morc\mtimescatmor \mord)}, \prreadout_{(\mora\mtimescatmor \morb)\mthen(\morc\mtimescatmor \mord)}, \prstart_\mora\tupconcat \prstart_\morb\tupconcat \prstart_\morc\tupconcat \prstart_\mord},
        \end{equation*}
    \end{widepar}
    with
    \begin{widepar}
        \begin{equation*}
            \defmapcomma{\prdyn_{(\mora\mtimescatmor \morb)\mthen(\morc\mtimescatmor \mord)}}
            {\prin_\mora\cprod \prin_\morb\cprod \prst_\mora \cprod \prst_\morb\cprod \prst_\morc \cprod \prst_\mord}
            {\sto}
            {\prst_\mora \cprod \prst_\morb\cprod \prst_\morc \cprod \prst_\mord}
            {\prinel_\mora \tupconcat \prinel_\morb\tupconcat \prstel_\mora \tupconcat \prstel_\morb \tupconcat \prstel_\morc \tupconcat \prstel_\mord}
            {\underbrace{\prdyn_{\mora \mtimescatmor \morb}(\prinel_\mora \tupconcat \prinel_\morb \tupconcat \prstel_\mora \tupconcat \prstel_\morb)
                    \tupconcat \prdyn_{\morc\mtimescatmor \mord}(\prreadout_{\mora \mtimescatmor\morb}(\prstel_\mora \tupconcat \prstel_\morb)\tupconcat \prstel_\morc \tupconcat \prstel_\mord)}_{(3)}}
        \end{equation*}
    \end{widepar}
    with
    \begin{widepar}
        \begin{equation*}
            \begin{aligned}
                (3) & =\prdyn_\mora(\prinel_\mora \tupconcat \prstel_\mora)\tupconcat \prdyn_\morb(\prinel_\morb \tupconcat \prstel_\morb)
                \tupconcat \prdyn_{\morc\mtimescatmor \mord}(\prreadout_\mora(\prstel_\mora)\tupconcat \prreadout_\morb(\prstel_\morb)\tupconcat \prstel_\morc \tupconcat \prstel_\mord) \\
                    & =\prdyn_\mora(\prinel_\mora \tupconcat \prstel_\mora)\tupconcat \underbrace{\prdyn_\morb(\prinel_\morb \tupconcat \prstel_\morb)}_{(*)}
                \tupconcat \underbrace{\prdyn_{\morc}(\prreadout_\mora(\prstel_\mora)\tupconcat \prstel_\morc)}_{(**)}\tupconcat \prdyn_\mord(\prreadout_\morb(\prstel_\morb)\tupconcat \prstel_\mord)
            \end{aligned}
        \end{equation*}
    \end{widepar}
    and
    \begin{equation*}
        \defmapcomma{\prreadout_{(\mora\mtimescatmor \morb)\mthen(\morc\mtimescatmor \mord)}}
        {\prst_\mora \cprod \prst_\morb\cprod \prst_\morc \cprod \prst_\mord}
        {\sto}
        {\prout_\morc\cprod \prout_\mord}
        {\prstel_\mora \tupconcat \prstel_\morb \tupconcat \prstel_\morc \tupconcat \prstel_\mord}
        {\prreadout_{\morc\mtimescatmor \mord}(\prstel_\morc \tupconcat \prstel_\mord)=\prreadout_\morc(\prstel_\morc)\tupconcat \prreadout_\mord(\prstel_\mord)}
    \end{equation*}
    As one can see from the expression for (3), the two terms ($*$) and ($**$) are switched compared to (1).
    Apart from this switch (and the corresponding switch in the signatures of the dynamics maps), we can see that there is a ``moral correspondence'' between the Moore machines in the functorial stacking axiom.
\end{example}




\subsection{Associative stacking}

\begin{ctdefinition}[Strict associative stacking semicategory]
    \label{def:strict-assoc-stacking-semicat}
    An \maindef{strict associative stacking category} is 
    
    \constit 
    \begin{enumerate}
        \item a \SY{functorial stacking semicategory} $\tup{\CatC, \mtimescat}$;
    \end{enumerate}
  
    \condit
    \begin{enumerate}
        \item the two composite functors $(( - ) \mtimescat ( - )) \mtimescat ( - )$ and $( - ) \mtimescat (( - ) \mtimescat ( - ))$ are equal as functors $\CatC \times \CatC \times \CatC \to \CatC$.
    \end{enumerate}
\end{ctdefinition}

\subsection{\Moore is associative stacking}
When considering Moore machines, we can define stacking operations and show that \Moore forms a stacking \SY{semicategory} (\cref{def:simple-stacking-semi-cat}).
The objects of \Moore are objects of \SetL, and therefore the stacking operation for objects corresponds to the ``multiplication in \SetL'', denoted by $\cprod$.

The operation on morphisms ``stacks'' Moore machines onto each other.
Formally:

\begin{equation}\label{eq:moore-stacking-prftree}
    \prfcomma{\mora\colon \prin_\mora \mtoin{\Moore} \prout_\mora}{\mora\colon \prin_\morb \mtoin{\Moore} \prout_\morb}
    {\mora \mtimescatmor \morb=\tup{\prin_\mora \cprod \prin_\morb, \prst_\mora \cprod \prst_\morb,\prout_\mora \cprod \prout_\morb, \prdyn_{\mora \mtimescatmor \morb},\prreadout_{\mora \mtimescatmor \morb}, \prstart_\mora \tupconcat \prstart_\morb}}
\end{equation}
with
\begin{equation*}\label{eq:moore-stacking-dyn}
    \defmapcomma{\prdyn_{\mora \mtimescatmor \morb}}{\prin_\mora \cprod \prin_\morb \cprod \prst_\mora \cprod \prst_\morb}
    {\sto}{\prst_\mora \cprod \prst_\morb}{\prinel_\mora \tupconcat \prinel_\morb \tupconcat \prstel_\mora \tupconcat \prstel_\morb}
    {\prdyn_\mora(\prinel_\mora \tupconcat \prstel_\mora)\tupconcat \prdyn_\morb(\prinel_\morb \tupconcat \prstel_\morb)}
\end{equation*}
and
\begin{equation*}\label{eq:moore-stacking-ro}
    \defmapperiod{\prreadout_{\mora \mtimescatmor \morb}}
    {\prst_\mora \cprod \prst_\morb}{\sto}{\prout_\mora \cprod \prout_\morb}
    {\prstel_\mora \tupconcat \prstel_\morb}{\prreadout(\prstel_\mora)\tupconcat \prreadout(\prstel_\morb)}
\end{equation*}
While we have already proved that the operation $\cprod$ is associative, it is also easy to see that the stacking of Moore machines is associative.
Therefore, \Moore equipped with the described stacking operations forms an \SY{associative stacking} \SY{semicategory}.




\subsection{Monoidal stacking}

\begin{ctdefinition}[Strict monoidal stacking semicategory]
    \label{def:strict-monoidal-stacking-semicat}
    A \maindef{strict monoidal stacking semicategory} is a \SY{stacking semicategory} $\tup{\CatC, \mtimescatob, \mtimescatmor}$ with

    \constit

    \begin{enumerate}
        \item an object $\idmoncat \setin \Obof{\CatC}$, called the \emph{monoidal unit}
    \end{enumerate}

    \condit

    \begin{enumerate}
        \item For any object $\Obja$ of \CatC,
              \begin{equation}\label{eq:strict-mon-stacking-semicat-cond1}
                  \Obja \mtimescatob \idmoncat = \Obja \qquad \text{and} \qquad \idmoncat \mtimescatob \Obja = \Obja.
              \end{equation}
        \item The monoidal unit $\idmoncat$ has an \SY{identity morphism} $\catid_\idmoncat$, and for any morphism $\mora \colon \Obja \mto \Objb$,
              \begin{equation}\label{eq:strict-mon-stacking-semicat-cond2}
                  \mora \mtimescatmor \catid_\idmoncat = \mora \qquad \text{and} \qquad \catid_\idmoncat \mtimescatmor  \mora = \mora.
              \end{equation}
    \end{enumerate}
\end{ctdefinition}

\todotext{J: check if the example below is in the right place, or if it should be a in a different subsection}

\begin{example}
    We can look at \Moore and ask whether it is a strict monoidal semicategory.
    The monoidal unit is given by the object
    $$\idmoncat = \cObj{}.
    $$
    Its identity morphism is the Moore machine
    \begin{equation*}
        \catid_\idmoncat=\tup{\cObj{},\cObj{},\cObj{},\prdyn_\idmoncat, \prreadout_\idmoncat, \tup{}},
    \end{equation*}
    where
    \begin{equation*}
        \defmapcomma{\prdyn_{\idmoncat}}
        {\cObj{} \cprod \cObj{}}
        {\sto}
        {\cObj{}}
        {\tup{}\tupconcat \tup{}}
        {\tup{}}
    \end{equation*}
    and
    \begin{equation*}
        \defmapperiod{\prreadout_{\idmoncat}}
        {\cObj{}}
        {\sto}
        {\cObj{}}
        {\tup{}}
        {\tup{}}
    \end{equation*}
    Clearly, $\prObja \cprod \cObj{}=\cObj{}\cprod \prObja=\prObja$ for every $\prObja\setin \Ob_\Moore$.
    Furthermore, consider a Moore machine $\mora\colon \prinL \mto \proutL$ with
    \begin{equation*}
        \mora=\tup{\prinL,\prstL, \proutL, \prdyn,\prreadout,\prstart}.
    \end{equation*}
    One has:
    \begin{equation*}
        \begin{aligned}
            \mora \mtimescatmor \catid_\idmoncat & =\tup{\prinL \cprod \cObj{},\prstL \cprod \cObj{},\proutL \cprod \cObj{},\prdyn_{\mora \mtimescatmor\catid_\idmoncat}, \prreadout_{\mora \mtimescatmor\catid_\idmoncat}, \prstart\tupconcat \tup{}} \\
                                                 & =\tup{\prinL,\prstL ,\proutL,\prdyn, \prreadout, \prstart}=\mora,
        \end{aligned}
    \end{equation*}
    where we used
    \begin{equation*}
        \begin{aligned}
            \prdyn_{\mora \mtimescatmor\catid_\idmoncat}\colon \prinL \cprod \cObj{}\cprod \prstL \cprod \cObj{} & \sto \prstL \cprod \cObj{} \\
            \prinel \tupconcat \tup{}\tupconcat \prstel \tupconcat \tup{}                                        & \mapsto \prdyn(\prinel,\prstel)\tupconcat \prdyn_\idmoncat(\tup{},\tup{})=\prdyn(\prinel,\prstel)
        \end{aligned}
    \end{equation*}
    and
    \begin{equation*}
        \begin{aligned}
            \prreadout_{\mora \mtimescatmor\catid_\idmoncat}\colon \prstL \cprod \cObj{} & \sto \proutL \cprod \cObj{} \\
            \prstel \tupconcat \tup{}                                                    & \mapsto \prreadout(\prstel)\tupconcat \prreadout_\idmoncat(\tup{})=\prreadout(\prstel)
        \end{aligned}
    \end{equation*}
    to show the equivalences $\prdyn=\prdyn_{\mora \mtimescatmor \catid_\idmoncat}$ and $\prreadout=\prreadout_{\mora \mtimescatmor \catid_\idmoncat}$.
    The argument for $\catid_\idmoncat \mtimescatmor \mora$ follows analogously.
\end{example}



% \begin{example}
\todo{Important: add comment regarding the fact that \Moore is not functorial stacking on the nose, therefore cannot be considered here?}
% \end{example}

\todotext{J: @JL: is Moore an example?}

\todotext{J: @J: what is a good definition of symmetric strict monoidal \SY{semicategory} ? Take inspiration from the \SY{adjunctions} for \SY{semicategories} paper that AC posted recently in order to encode ``isomorphism'' without saying it with identities?}

